\documentclass{article}
\usepackage{mathnotes}

\title{Classification of Singularities}
\description{Types of isolated singularities including removable singularities, poles of various orders, and essential singularities based on Laurent series.}

\begin{document}

\begin{definition}[Singularity]
A point $z_0$ is called a \textbf{singularity} of a complex function $f$ if $f$ is not analytic at $z_0$, but every neighborhood of $z_0$ contains at least one point at which $f$ is analytic.
\end{definition}

\begin{definition}[Point Singularity]\synonyms{"Isolated Singularity"}
A \dref{singularity} $z_0$ of a \dref{complex-function} $f$ is said to be a \textbf{point singularity} or \textbf{isolated singularity} if there exists a \dref{neighborhood} of $z_0$ in which $z_0$ is the only \dref{singularity} of $f$.

Alternatively, a \textbf{point singularity} of a function $f$ is a \dref{complex-number} $z_0$ such that $f$ is defined in a \dref{neighborhood} of $z_0$ but not at the point $z_0$ itself.
\end{definition}

\section{Classification of Singularities}

Given a \dref{Laurent-series} expansion of a \dref{function} with an \dref{isolated-singularity} at $z_0$, we have the possibility of the \dref{singularity} being removable, a pole, or an essential singularity.

\begin{definition}[Removable Singularity]\synonyms{removable}
A \textbf{removable singularity} occurs if all negative powers in the Laurent series are zero and the function can be redefined at the singularity to be analytic.
\end{definition}

\begin{example}
Consider $f(z) = z$ defined in the \dref{punctured-plane}. Then, the \dref{origin} is an \dref{isolated-singularity}, because $f(z)$ is not defined there, but is defined everywhere else in a \dref{neighborhood} of the origin. We can easily define $f(0) = 0,$ and this will make $f(z)$ analytic at $0,$ and so we say this \dref{singularity} is \dref{removable}.
\end{example}

\begin{definition}[Pole]
A \textbf{pole} is present if the Laurent series has a finite number of negative power terms. The largest negative exponent (in absolute value) indicates the order of the pole.
\end{definition}

\begin{example}
Consider $f(z) = 1/z$ defined in the \dref{punctured-plane}. Then, the \dref{origin} is an \dref{isolated-singularity}, but in this case, we can't simply define it in a \dref{continuous} fashion, because $f(z)$ approaches infinity as $z \to 0.$ Thus, we call this \dref{singularity} a \dref{pole}.
\end{example}

\begin{definition}[Essential Singularity]
An \textbf{essential singularity} occurs if there are infinitely many negative powers in the Laurent series.
\end{definition}

\begin{example}
Consider $f(z) = e^{1/z},$ on the \dref{punctured-plane}. As $z$ approach $0$ on the positive real axis, $f(z)$ goes to infinity, and as $z$ approaches $0$ on the negative real axis, $f(z)$ approaches 0. On the imaginary axis, $f(z)$ oscillates wildly, but remains \dref{bounded}, as $z$ approaches $0.$ This is neither a \dref{removable-singularity} nor a \dref{pole}, and we call it an \dref{essential-singularity}.
\end{example}

\begin{theorem}
If $f(1/z)$ has a \dref{singularity} at $0,$ then $f(z)$ has a \dref{singularity} at the \dref{point-at-infinity}.
\end{theorem}

\end{document}
